% !TEX root = ../main.tex
% !TEX program = xelatex
\title{融合多源预测的微网经济调度与线性优化研究}
\date{}
\maketitle
\begingroup
\setstretch{1.22}
\begin{abstract}
\noindent\hspace*{2em}随着我国碳达峰、碳中和目标的推进和分布式新能源的发展，光伏与储能逐步成为社区供能的重要组成。如何在供需和电价变化下合理购电，既避免高价紧急补购，又减少多订电量的损失，是微网运行的关键。本文以10分钟为基本时段，以电量平衡和储能约束为共同基础，逐步引入历史预测、日内预报和电价预测，形成由确定性调度向动态调整递进的购电策略。

\textbf{针对问题一}，采用\textbf{电量流线性规划}，统筹全天购电及储能充放电，在保持日初、日末储电量相等的条件下，将低价购电和富余光伏留到更有价值的时段使用。当天最低购电费为\textbf{35126.85元}，较峰谷策略降低\textbf{8.23\%}，体现了供需与储能联合安排的节费优势。

\textbf{针对问题二}，将\textbf{历史分项预测、近期加权准备与跨日线性规划}相结合。负载预测利用星期规律，光伏预测参考近期发电水平，再根据近期缺口误差设置购电准备，并通过48小时规划衔接次日用电。2025年2—12月总费用为\textbf{1358.38万元}，较等权误差准备节省\textbf{44.66万元}。结果表明，适量增加普通购电准备能够减少更昂贵的紧急补购，预测与调度的配合比单独追求预测精度更贴合本题目标。

\textbf{针对问题三}，运用\textbf{融合预测与日内校正滚动规划}，将官方光伏预报与近期历史结合，利用当天已发生的负载偏差修正后续判断，并在6、12、18点权衡调整费与补购费，更新剩余订单。同期总费用为\textbf{1301.42万元}。同一模型下的信息对照表明，日内反馈与改单节省\textbf{50.63万元}；在此基础上，后发官方预报再节省\textbf{6.95万元}，并使紧急购电量下降\textbf{30.10\%}。两组期末库存相同，说明日内新预报能够带来额外节费，并减少对紧急补购的依赖。

\textbf{针对问题四}，在未来真实电价未知的条件下，零点分支使用\textbf{提升树预测与线性规划}，利用历史价格和供需信息刻画价格变化；日内分支选取\textbf{岭回归预测与滚动线性规划}，在候选比较中以接近最低的费用获得较少的紧急购电，并便于解释价格变化的影响。两套完整策略同期费用分别为\textbf{1424.86万元、1369.59万元}，日内策略节省\textbf{3.88\%}，体现了供需判断、价格预测与订单调整的综合收益。

逐10分钟复核表明，各方案满足供需平衡和储能设备约束，费用可由实际订单逐项复算。参数比较揭示了提前准备与高价补购之间的费用权衡，信息对照进一步区分了反馈调整和新增预报的贡献。本文以\textbf{实际购电费用}连接预测与调度，为社区微网兼顾用能经济性与可再生能源利用提供了可核验的方案。
\end{abstract}
\noindent{\heiti\zihao{-4} 关键词：}{\songti\zihao{-4}\textbf{微网调度；线性规划；误差准备；滚动优化；电价预测}}
\par
\endgroup
